\documentclass[]{../../../../tex/classes/styledArticle}
\usepackage{../../../../tex/packages/hebrewSupport}


\author{שחר פרץ}
\title{\textit{היסטוריה 52} – מסיבת ציוני הגשה}
\begin{document}
	\maketitle
	ההערות שנעבור עליהם עכשיו בנוגע לבחינה, לא רק שלנו. אלא הערות כלליות ודברים ששרית שמעה ממורים בצוות. נתחיל מהערות כלליות, ונעבור להערות לכל שאלה בנפרד. 
	\begin{itemize}
		\item כאשר מבקשים קטע מקור ועובדה היסטורית (נאמר, שתי עובדות ומקור, או עובדה ומקור) – אל תקחו את אותה העובדה ומקור. זה אותו הדבר. לא יקבלו את זה. נוסף על זה, אם טעיתם טעיתם פעמיים. 
		
		דוגמה: היה כתוב ''בססו על שתי עובדות היסטוריות ומקור מספר הלימוד`` – אם העובדה לא נכונה, ואתם מבססים את אותה העובדה על מקור היסטורי, איבדתם 2/3 מהשאלה. 
		\item קטע מקור – אל רק תצטטו ותסבירו במילים שלכם. צריך לקשר. 
		\item אפשר להביא מקור תמונה. קשה יותר לבסס את זה. צריך שני רכיבים לבסס את הקישור של המקור, מתמונה יותר קשה. לפעמים יש טקסט ותמונה צמודים, ואתם בוחרים את התמונה. 
		\item תתכננו את התשובות. יש כאן לא מעט ילדים שעשו קווים. כשאתם עושים דברים כאלו מעריך יכול לפספס. 
		\item חלק מכם בפתיחה מציגים יפה את המקור שאתם הולכים להביא מהספר, ושניים וחצי עמודים אחכ אתם מתייחסים אליו. 
		\item ארץ זו לא מדינה. ארץ זה אזור גיאוגרפי, מדינה זה יישוב פוליטי. 
		\item שמות מלאים. אריאל שרון לא היה חבר שלכם. אל תכתבו ''אריאל אמר``. אפשר לכתוב שם משפחה. גם לא לכתוב ''מנחם`` (בגין). לא יורידו על זה ניקוד אבל זה לא נראה טוב. 
		\item חלקכם בכלל לא מתייחסים למקור, אלא רק להבדלים וכו'. 
		\item כשיש טענה/טיעון/עמדה – צריך להציג את הטיעון. 
		\item סיבות להבדלים בין קטעי מקור מתייחסים לפרטי המקור. כאשר מדברים על הבדלים ביחס בשאלה 1, מדברים על הבדלים בין קטעי המקור. 
	\end{itemize}
	ברמה הטכנית: 
	\begin{itemize}
		\item תמרקרו
		\item מחברות הבגרות הן מחברות מתמטיקה, שבין כל שתי שורות יש פס שאף אחד לא רואה, שאמור לסמן את השורות. כאשר אתם כותבים בתוך המשבצות, ''זה הסיוט של כל מעריך``, אי אפשר לקרוא את זה. ואל תכתבו בשני הצדדים. 
	\end{itemize}
	
	עכשיו נתמקד בשאלה שאלה. 
	\begin{enumerate}
		\item ראשונה – ''בתהליך בניית המשטר הנאצי הנאצים השתמשו בשלושה תחומים... בחרו שניים, בחרו מקור ועובדה היסטורית שמבטאים כיצד הם גייסו את תמיכת ההמונים``. 
		\begin{itemize}
			\item זו לא שאלה על יהודים והיא לא מכוונת לכך. אפשר לקשר את זה עם ממש רוצים. 
			\item מוקש ראשון – הניסוח הוא תהליך בניית המשטר הנאצי. הוא הסתיים באוגוסט 34 כשהנשיא הידנבורג מת. ``ליל הבדולח כתבת?\ \!``
			\item מוקש שני – הנושא של גיוס ההמונים. טרור ותעמולה בדכ די קל להסביר. חקיקה פחות. אפשר חקיקה, אבל צריך להבחור חוק שקל להסביר את הנושא של גיוס ההמונים. היה מי שכתב על חוק ההסמכה, הסביר את החוק ואז קבע בלי קישור שהוא גייס המונים. 
		\end{itemize}
		
		עכשיו נדבר על הסעיף הבא. 
		\begin{itemize}
			\item ''זה פשוט מדהים להבין כמה אתם לא מבינים מה זה הסדר החדש``. היה צריך להסביר כיצד עקרונות (ברבים) באים לידי ביטוי בסדר החדש. 
			\item כאשר אתם כותבים על תורת הגזע, זה נכון שיש בה מדרג גזעי וכו', אבל לא הגדרתם. תזכרו ''הדם הוא האדם`` – חלוקה לגזעים, כאשר לכל גזע יש תכונות פנימיות וחיצוניות שלא ניתנות לשינוי. אחכ הם מחלקים את העולם לפירמידת גזעים. 
			\item מרחב מחייה – בשטחי המזרח
			\item אנטישמיות היא גזעית, אבל זה לא בא ביחד עם מרחב מחיה. 
		\end{itemize}
		אחכ ביקשו להציג את המדיניות ליהודים עפ כל אחד מקטעי המקור. 
		\begin{itemize}
			\item יש מדיניות ויש צעדים. המדיניות היא ''בידוד``, הצעדים הם ''תלאי צהוב, הוצאה ממערכת החינוך``, וכיו''ב. 
			\item הרבה פעמים יש קשר בין הקווקווים בסעיפים. היה קווקוו של להסביר על הסדר החדש, ולאחר מכן קווקוו של למה יש הבדל. תחשבו, אולי יש קשר. 
			
			יש ילדים שכתבו שההסבר הוא בזמנים – אין קשר. כל טקסט נכתב בשלב בו המדינה נכבשה. המדיניות נקבעת לפי המדרג הגזעי בסדר החדש. הדוגמה הפשוטה – כאשר לא־יהודים הצילו יהודים בפולין, לעומת הולנד, הענישה הייתה שונה לחלוטין. 
		\end{itemize}
		\item \begin{enumerate}[A']
			\item בשנים 39-40 ערער המעבר לגטו את חיי [...] הילדים בפרט. הסבירו טענה זו ע''פ מקור בספר הלימוד. 
			\begin{itemize}
				\item התייחסו לילדים. יש מי שהתחיל לפרוס את קשיי היהודים בגטו. זה לא קשור באופן ישיר, צריך להסביר את הקשר. 
				\item הרעיון די ברור – הילדים שנאלצים להיות מפרנסים, היפוך תפקידים. כתבתם על התמונה של הילדים שמבריחים – איך זה מבטא את ערעור החיים? מה המשמעות של זה שילד צריך להביא אוכל? שמערכות חינוך לא קיימות? 
			\end{itemize}
			ישנו חלק שני לסעיף זה. 
			בסעיף השני – בברהמ חל שינוי בברית המועצמות. הסבירו את השינוי. 
			\begin{itemize}
				\item כאשר יש שינוי, היה, והשתנה. צריך להסביר מה היה. 
				\item לא לכוון לדעת הפירהרר. זו תזה איזוטרית. 
				\item בספר יש הסבר לא משהו לגבי זה שלנאצים לא היו תוכניות הגירה בכלל כמות היהודים. זה בעייתי. היו שתי תוכניות להגירת יהודים (מדסקר, ואוניסקו־לובלין), שתיהן הסתיימו עד קיץ 1940. הן לא רלוונטיות, זה ממש טעות. מצד שני, מי שכתב את זה כתב על כמות היהודים, ואז הוא קיבל על זה נקודות. 
				\item יש הסברים שכתובים בספר ובסיכום בצורה ברורה. אל תשתמשו בהסברים איזוטריים. 
			\end{itemize}
			\item חסידי אומות עולם. סיפורם של הזוג ז'בלינסקי או משהו כזה. הסבירו למה הם חסידי אומות עולם, והסבירו את ההתאמה. 
			\begin{itemize}
				\item צריך אשכרה להסביר את התנאים. אחד מהם, אומר שיהיה להם תיעוד. בסופו של יום לקחו רק שני תנאים, בבגרות ידרשו שלושה. 
				\item דפוסים – חסידי אומות עולם זה לא דפוס בפני עצמו. זה חלק מהדפוס של הצלה. יש הצלה ע''י חסידי אומות עולם ויש הצלה שאינה. 
				\item שנה שעברה מישהו כתב שז'בלינסקי הוא יהודי כי יש לו שם יהודי. שמות יהודים ופולנים נשמעים דומים. 
				\item דפוס יחס אחר – להביא את היחס. בקשו גם מקור – ממש עדיף טקסט. כאשר יש תמונה של אנשים צופים בנאצים יורים על יהודים לתוך הבורות, איך אתם מוצאים שני מאפיינים? 
			\end{itemize}
		\end{enumerate}
		\item \begin{enumerate}[A']
			\item 
			\begin{itemize}
				\item קראו את קטע המקור וענו על השאלה... היה קטע מקור על כיבוש צרפת. בקשו ע''פ קטע המקור ולפי מה שלמדנו. אחד כך אחד כך. עשיתם שלושה – השלישי יכול להיות גם וגם. בקטע הזה הבליצקריג זה הדבר המובהק. בקשו סיבה. צריך סיבה באופן כללי, לא עמוד שלם על צרפת ולשכוח מה הסיבה שהם נצחו באופן כללי בשנתיים הראשונות. 
				\item החל משנת 1942 חל שינוי במצבה של גרמניה – \textit{שינוי}. צריך להסביר מה היה קודם, ומה השתנה. בקשו \textit{שינוי במלחמה} – הפתרון הסופי לא רלוונטי פה. ההסבר לא כזה מסובך – יש לנו קרבות לפני. רק צריך להסביר שבשנתיים הראשונות גרמניה מנצחת. 
				\item היה צריך לבחור באחד הקרבות. זה היה 42-44, אז כיבוש ברלין לא רלוונטי. 
				\item הסבירו את הקרב והסבירו במה הוא תרם לשינוי. בלי התרומה, זו לא תשובה. המהות של השאלה היא לקחת דוגמה ולהראות איך היא תרמה לשינוי. 
			\end{itemize}
			\item היסטוריונים טוענים ששואת יהודי צפון אפריקה היא חלק משואת היהודים. \begin{itemize}
				\item הציגו 2 סיבות להבדל, בססו את הטענה, וכו'. 
				\item אמרו טענה – תבססו את הטענה. 
				\item כאשר מציגים את פרוטוקול ועידת ואנזה, צריך להסביר איפה הם נמצאים בפרוטוקול. לא רק לציין שהם הופיעו. ממש טבלה, מספר, וכו'. 
			\end{itemize}
			overall שאלה ארוכה אבל מי שענה ידע את התשובה טוב. 
		\end{enumerate}
		
		המורה אמרה מאוד בזהירות (עניין אידבדואלי) – סטטיסטית, מי שענה על שאלה 3, ירדו לו פחות נקודות. מצד שני, יש ילדים שיבחרו את שאלה 1, שהייתה מצוינת בשבילם, למרות שסטטיסטית ירדו הכי הרבה נקודות בה. 
		\item \begin{enumerate}[A']
			\item גורמים מסייעים ומעכבים, 2 מטרות ו־2 דרכים לאחת המטרות. \begin{itemize}
				\item שואלים על מטרות. זה ברבים. 
				\item דרכי פעולה 
			\end{itemize}
			\item יש איזה קטע על רוטשילד כאן. שרית עשתה הדגמה על טקסט של עמדה על רוטשילד, הפוכה – של מישהו שטוען שהשיטה טעות, ומבקשים להציג עמדה שונה. סתם טעות במצגת. 
			\begin{itemize}
				\item בקשו להסביר ע''פ הקטע את עמדת ההיסטוריון. צריך להציג את עמדת ההיסטוריון. לא לכתוב ''עמדת ההסטריון חיובית`` – מה זה חיובית? 
				\item בססו על שתי עובדות מהמקור – זה אומר שצריך לבסס על שתי עובדות, לא לקחת ציטוט ולהסביר אותו במילים שלנו. 
				\item הביאו מקור מספר הלימוד שמבטא עמדה שונה – היא מקור אחד בלבד. כולם מצאו אותו. הבעיה היא שהוא לא כותב דברים ככ מפורשים, אבל הוא מדבר על היהודי החדש, ועל זה שרוטשילד מנוגד אידיאולוגית לכיבוש הקרקע והעבודה. הוא פשוט אומר את זה אחרת. 
			\end{itemize}
		\end{enumerate}
		
		\item \begin{enumerate}[A']
			\item כולם ענו טוב. צריך לשים לב שבקשו שני שלבים של מאבק, ומנהיג שהשפיע על שלבי המאבק. שימו לב שקודם תוודאו שהשלבים שאתם מדברים זה מה שתוכלו לקשר אחכ למנהיג. 
			\item יש הטוענים שהרצל הוא מנהיג מדינת היהודים הודות לחזון ולכישוריו הדיפלומטיים. רעיון=חזון. כינוס הקונגרס הוא לא דיפלומטיה. זה גם לא חזון. חזון זה רעיון. אפשר את תוכנית באזל. מי שכתב את כינוס הקונגרס כדיפלומטיה (אני), הפסיד את כל הניקוד. 
		\end{enumerate}
		
		\item אתגר זה לא גורם
		\item \begin{enumerate}[A']
			\item השאלה עם המדליות. בקשו קשיים \textit{במבצע ההעפלה}. לא קשיים של המעפילים. לכתוב שהיה צפוף באוניות, לא קשור. 
			\item חלק ענו כאן מהמם וחלק לא הבינו מה רוצים מהחיים שלהם. ההיסטוריון יובל נח הררי טוען שמי שטוען שפעילות המחתרות גרמה ליציאת בריטניה, טועה ומטעה. 
			\begin{itemize}
				\item לטענתו, אין קשר למחתרות. הן לא הרסו, לא עזרו. אין קשר. העולם לא מסתובב רק סביבנו, קוראים עוד כמה דברים. כל מי שענה משהו שקשור למחתרות, לדוגמה, הטיעונים ללמה אירועים מסויימים פגעו ביחס של דעת הקהל העולמית, לא קשור. 
				\item לא חסרים עובדות היסטוריות שלא קשורות למחתרות – עליהם צריך לכתוב. 
			\end{itemize}
		\end{enumerate}
		\item ''מורכבת אך מאוד קלה למי שלמד מלחמת העצמאות`` \begin{enumerate}[A']
			\item \begin{itemize}
				\item הציגו את החלטת האו''ם – שתי מדינות לשני עמים. אם זה לא כתוב מורידים נקודות. 
				\item שימו לב, ביקשו על הקשיים שנובעים מההחלטה. הקשיים בהחלטה הם הקשיים של השלב הראשון במלחמה. שנבעו מההחלטה באופן ישיר. 
				\item הציגו שני סעיפים (בהחלטה) שלא מומשו בסיום המלחמה. ''בתכלאס, אף סעיף לא מומש``. שיתוף פעולה כלכלי בין שתי המדינות אומנם לא מומש, אבל ''נו באמת``, יותר הגיוני לדבר על ירושלים הביןלאומית, על הגבולות, וכו'. 
			\end{itemize}
		\end{enumerate}
		\item יש מישהו אחד שענה על 9. 
		\begin{enumerate}[A']
			\item 
			\item משהו לגבי זה שמנהלת העם הייתה מחלוקת לעיתוי ההחלטה. כאשר דיברו על ידין, אמרו מפורשות שהוא בא מרקע צבאי. הוא מדבר על הכוחות האוויריים שלהם, ושהסיכויים שקולים לשני הצדדים. ''אם להיות יותר כנה, הייתי אומר שהייתרון שלהם גדול``. ביקשו להסביר את הטענה ולתת טיעון נגדי. יש את זה פשוט בספר. 
		\end{enumerate}
		\item דה־קולוניזציה – גם אם זה לא נדרש, תגדירו זמן מקום וכו', זה יעזור לכם. 
		\begin{enumerate}[A']
			\item מלחמה קרה באה בחלק הפוליטי ובחלק האידיאולוגי. כנל על הנושא הלאומי. זה הסברים שונים. 
			\item שאלו על היהודים, לא הייתה בעיה. נקודה אחת – פאן־ערביות זה לא דתי, אלא לאומי. 
		\end{enumerate}
		\item לא היה
		\item לא היה
		\item \begin{enumerate}[A']
			\item הציגו שלוש מהמטרות [...] אילו מהמטרות הושגו בתום המלחמה. 
			\begin{itemize}
				\item החזרת שטחים לא הושגה. 
				\item הסכם שלום לא הייתה מטרה. השגת תהליך מדיני כן. וזה לא נכון שאחרי הסכם המלחמה החל תהליך מדיני. אומנם היו הסכמי ביניים, אבל ההסכמיים המדיניים מתחילים מאוחר יותר. נסללה הדרך להסכמים. 
				\item השאיפה לשקם את הכבוד – תסבירו למה. לכתוב שזה ככה כי ''המנטליות הערבית מבוססת על כבוד``, וכל מני הסברים על המנטליות הערבית, לא הולך. 
			\end{itemize}
			\item \begin{itemize}
				\item בקשו סיבה משותפת לשני הטקסטים. חלקכם חירטטתם. אין קשר לקונספציה. יש פה זלזול באויב, ואפשר ואף כדאי להסביר שהקונספציה גרמה לזלזול באויב. 
				\item הסבירו כיצד עובדה זו השפיעה על הימים הראשונים. אנשים לא ענו על זה, או שהם ענו ''היו מלא הרוגים`` (לא מלא מספיק). 
				\item בקשו הבדל בין שני קטעי המקור בתיאור כוחו ומעמדו של צה''ל, ושתי סיבות להבדל. הן תלויות במי, מתי ואיפה. 
				
				לא ירד על זה נקודות למרות שזה כתוב במחוון, אבל שרון כבר היה בדרך לפוליטיקה ובשלב הזה היה עם אג'נדה. בסופו של יום יש לו ראיון בעיתון מעריב. גם שרון לחם במלחמה. יש גם את העניין של בדיעבד. 40 שנה אחרי, יש לזה משמעות. זה לא סתם 40 שנה אחרי, יש כאן פרספקטיבה אחרת, הסקת מסקנות, וכו'. 
			\end{itemize}
		\end{enumerate}
	\end{enumerate}
	
	דרוזים אין במבחן הזה. אבל שימו לב – קרב רמת יוחנן לא דוגמה לסיוע של הדרוזים. מבצע חריה. 
	
	היום ב־12:45 עד 13:30. תלמידים שרוצים לבוא, לערער, ששרית תבוא ותסביר, מוזמנים לבוא. מי שלא הספיק לדבר, שרית תמצא זמן להסביר לו. מרכז משאבים. ציוני ההגשה מורכבים מכל הפריסה. 30\% מתכונת. אם לילד יש ציון חריג (לכאן ולכאן), נתחשב בכך. ציוני שנה שעברה הוכנסו לציונים שותפים כדי שלא תצטרכו לחפש. מחר לשיעור אתם מגיעים עם החישוב הזה. מחר כולם יוצאים מחוץ לכיתה ואז אחד אחד תקבלו את ציוני ההגשה. 
	
\end{document}